% Auto-generated by scripts/python/bundle_paper_tex.py — do not edit by hand.
% Edit mub6_karlsson_ieee.tex and paper/proofs/*.tex, then regenerate.
% Overleaf: upload this file as main.tex (default) — no proofs/ folder needed.
% Local multi-file compile: use main_multifile.tex or mub6_karlsson_ieee.tex + proofs/.

% Technical report: MUB extendability within Karlsson K_6^(3) only.
% Intended for arXiv (quant-ph) + reproducible code repository.
%
% Compile from paper/ directory (required for proofs/*.tex \input files):
%   cd paper && pdflatex mub6_karlsson_ieee.tex && pdflatex mub6_karlsson_ieee.tex
% Overleaf / single-file: use main_standalone.tex (python scripts/python/bundle_paper_tex.py)
% Cursor/VS Code: .vscode/settings.json sets useFileDir=true so Build works from the IDE.

\documentclass[onecolumn]{IEEEtran}

\usepackage[T1]{fontenc}
\usepackage[utf8]{inputenc}
\usepackage{amsmath,amssymb,amsfonts}
\usepackage{graphicx}
\usepackage{cite}
\usepackage{booktabs}
\usepackage{array}
\usepackage{hyperref}
\usepackage{url}
\usepackage{xurl}
\usepackage{amsthm}
\usepackage[strings]{underscore}

% Allow line breaks in long repository paths and script names
\newcommand{\repo}[1]{{\ttfamily\nolinkurl{#1}}}
\emergencystretch=2.5em
\sloppy

\newtheorem{theorem}{Theorem}
\newtheorem{claim}{Claim}
\newtheorem{lemma}{Lemma}
\newtheorem{proposition}{Proposition}
\newtheorem{conjecture}{Conjecture}
\newtheorem{remark}{Remark}

\hypersetup{
    colorlinks=true,
    linkcolor=black,
    citecolor=black,
    urlcolor=blue
}

% Avoid bad line breaks in citations
\def\BibTeX{{\rm B\kern-.05em{\sc i\kern-.025em b}\kern-.08em
    T\kern-.1667em\lower.7ex\hbox{E}\kern-.125emX}}

\begin{document}

\title{Third-MUB Loci and Fourth-MUB Negative Evidence Within Karlsson's Three-Parameter Family in Dimension Six}

\author{%
\IEEEauthorblockN{Adian Jaupi}%
}

\maketitle

\begin{abstract}
Within Karlsson's three-parameter family $K_6^{(3)}$ only---not all CHMs in $\mathbb{C}^6$---I report proved gauge structure, HP-supported third-MUB evidence on the Dita $\lambda$-circle, numerical fourth-MUB negative evidence on that slice, and a reproducible certified pipeline.
\textbf{(T1) Gauge structure (proved):} at $\theta{=}0$, $\phi$ is a parametrization gauge; on the Dita slice ($\theta{=}\arccos(1/\sqrt3)$, $\phi{=}\pi/4$), distinct $\lambda$ yield CHM-inequivalent matrices (Theorem~\ref{thm:T1}, \S\ref{sec:proofs}).
\textbf{(T2) Third MUB on Dita $\lambda$-circle (HP-supported):} at each of $126/126$ periodicity samples and $628/628$ dense probes on the circle, $\{\mathbf{I},H(\lambda)\}$ admits a third MUB via certified pool extraction; extension to all $\lambda\in\mathbb{R}$ additionally uses $2\pi$-periodicity and an informal continuity argument (\S\ref{sec:proof-t2}).
\textbf{(T3) No fourth MUB on the Dita circle, confirmed numerically over $628/628$ probes;} a symbolic elimination witness ($35$ equations, $20$ variables) was constructed but computed over $\mathbb{C}$ with floating-point coefficients and without exact-arithmetic elimination, so it does not constitute a proof (\S\ref{sec:proof-t3}).
\textbf{(C1) Literature correction:} Karlsson's original $A$-matrix passes $348/349$ tests; McNulty--Weigert transcription fails $349/349$.
\textbf{(Numerical classification, C2/C3)} HP audit within the certified search CSV ($928$ rows, $899$ pool-complete) and Track~D classification identifies two reproducing third-MUB loci (F6 $\theta{=}0$ arc; Dita circle)---not a global family classification.
\textbf{($\mathcal{S}^*$ primary sweep)} $1863/1865$ distinct candidate points checked ($0$ fourth MUB); clique-6 counts are per candidate locus, not distinct third-basis solutions (\S\ref{sec:trackc}).
This report does not prove Zauner's conjecture $N(6){=}3$ or fourth-MUB absence on all of $K_6^{(3)}$.
Cross-family extension (Phase~C / Liang--Chen sources) is blocked pending source transcription and is not claimed here.
\end{abstract}

\begin{IEEEkeywords}
Complex Hadamard matrices, Karlsson family, mutually unbiased bases, numerical algebraic geometry, homotopy continuation, reproducible computation
\end{IEEEkeywords}

\section{Introduction}
\label{sec:intro}

Mutually unbiased bases (MUBs) in $\mathbb{C}^d$ are orthonormal bases $\mathcal{B}_0,\ldots,\mathcal{B}_{m-1}$ such that for all $a\neq b$ and all $|\psi\rangle\in\mathcal{B}_a$, $|\phi\rangle\in\mathcal{B}_b$,
\begin{equation}
    \label{eq:mu}
    \bigl|\langle\psi|\phi\rangle\bigr|^2 = \frac{1}{d}.
\end{equation}
MUBs play a central role in quantum state tomography, cryptography, and the study of finite geometry~\cite{Bengtsson2007, Brierley2009}.
For prime-power dimensions $d=p^n$, exactly $d+1$ MUBs exist~\cite{Ivanovic1981, Wootters1989}.
For composite $d$, the maximum number $N(d)$ of MUBs is unknown in general; in dimension six, Zauner conjectured $N(6)=3$~\cite{Zauner1999, Scott2010}.

A standard construction route passes through complex Hadamard matrices (CHMs): unimodular $d\times d$ matrices $H$ with $HH^\dagger=dI$.
Columns of $H/\sqrt{d}$ (after suitable normalization) define a basis mutually unbiased to the computational basis.
Finding additional MUBs amounts to finding further orthonormal bases whose columns are unbiased to existing ones.

Karlsson~\cite{Karlsson2011} parametrized a three-dimensional family $K_6(\theta,\phi,\lambda)$ of $6\times6$ CHMs, subsuming earlier two-parameter families due to Dita and others~\cite{Dita2009, McNulty2024}.
McNulty and Weigert~\cite{McNulty2024} survey CHM classification and conjecture (Conj.~8.1) that triples of MUBs arising from Karlsson matrices are rare or absent for generic parameters.
Grassl~\cite{Grassl2004} showed that the pair $\{\mathbf{I},F_6\}$, where $F_6$ is the cyclic Fourier matrix, extends to at most three MUBs---so searches must be \emph{per} Hadamard matrix, not via a fixed pool unbiased only to $F_6$.
Zauner's conjecture $N(6)=3$~\cite{Zauner1999, Scott2010} motivates the broader question but is \textbf{not} addressed globally here.

\subsection{Scope and Position Relative to Prior Work}
\label{sec:prior}
\textbf{Scope.} All results concern Karlsson's family $K_6(\theta,\phi,\lambda)$ only.
CHMs outside $K_6^{(3)}$---Bj\"orck circulants, other Di\c{t}a-type families, and constructions surveyed in~\cite{McNulty2024}---are not searched.
\textbf{Relation to prior work.}
Grassl~\cite{Grassl2004} bounds extendability from $\{\mathbf{I},F_6\}$; Brierley--Weigert~\cite{Brierley2009} established third MUBs at Dita-type Karlsson parameters; I extend this to a full $\lambda$-circle of CHM-inequivalent matrices (Theorem~\ref{thm:trackd}).
McNulty--Weigert Conjecture~8.1~\cite{McNulty2024} predicts rarity of third-MUB triples in Karlsson's family; my fourth-MUB numerics on the certified search CSV, $\mathcal{S}^*$ primary sweep ($1863/1865$), and dense Dita probes are \emph{confirmatory} heuristic evidence, not a proof.
The per-$H$ pipeline applies standard homotopy continuation~\cite{HomotopyContinuation}; its value lies in certified problem-specific workflow (corrected completeness gate, HP re-verification, formal search sets), not in new numerical algebraic geometry algorithms.

\subsection{Contributions Within \texorpdfstring{$K_6^{(3)}$}{K6(3)}}
This technical report documents:
\begin{enumerate}
    \item[\textbf{C1}] A \textbf{literature correction}: Karlsson's original $A$-matrix vs.\ the McNulty--Weigert transcription (Table~\ref{tab:audit}).
    \item[\textbf{C2}] The \textbf{main structural result}: exactly two HP-verified third-MUB loci---F6 $\theta{=}0$ bounded $\lambda$-arc and Dita full $\lambda$-circle---with CHM-inequivalence and contrasting topology (Theorem~\ref{thm:trackd}).
    \item[\textbf{C3}] \textbf{Confirmatory negative numerics}: no fourth MUB on the certified search CSV ($928$ rows), $\mathcal{S}^*$ primary ($1863/1865$ candidate points), dense Dita probes ($628/628$), or $\theta{=}0$ slices; consistent with Conj.~8.1, not a theorem on $\mathcal{R}$ or $N(6)$.
    \item[\textbf{C4}] A \textbf{reproducible certified workflow} (per-$H$ MU pool, clique extension test, search-set formalism; Appendix~\ref{app:reproduce}), including documented correction of an initial coarse-grid overclaim.
\end{enumerate}
The investigation used three internal tracks: classification (Track~D), fourth-MUB numerics on $\mathcal{R}\cap\mathcal{S}_{588}$ (Track~C), and dense probes on known third-MUB loci (Track~A).
A symbolic Groebner ladder (Section~\ref{sec:symbolic}) is exported but \textbf{incomplete}; this report is submitted to document C1--C4 while elimination remains future work.

\section{Background}
\label{sec:background}

\subsection{MUBs From Hadamard Matrices}
Let $H$ be a $6\times6$ CHM with $|H_{ij}|=1$ and $HH^\dagger=6I$.
The normalized columns $c_k=H_{:,k}/\sqrt{6}$ form an orthonormal basis $\mathcal{B}_H$ satisfying~\eqref{eq:mu} with the computational basis $\mathcal{B}_I=\{e_0,\ldots,e_5\}$.
Given two such bases, a third MUB $\mathcal{B}_3$ requires six unit vectors, pairwise orthogonal and unbiased to both $\mathcal{B}_I$ and $\mathcal{B}_H$; a fourth MUB requires a second disjoint six-vector set with the same property relative to $\mathcal{B}_3$.

\subsection{Karlsson's Three-Parameter Family}
Following Karlsson~\cite{Karlsson2011} and the review of McNulty and Weigert~\cite{McNulty2024}, write
\begin{equation}
    K_6(\theta,\phi,\lambda)=
    \begin{bmatrix}
        F_2 & Z_1 & Z_2 \\
        Z_3 & \tfrac{1}{2}Z_3 A Z_1 & \tfrac{1}{2}Z_3 B Z_2 \\
        Z_4 & \tfrac{1}{2}Z_4 B Z_1 & \tfrac{1}{2}Z_4 A Z_2
    \end{bmatrix},
\end{equation}
where $F_2=\bigl[\begin{smallmatrix}1&1\\1&-1\end{smallmatrix}\bigr]$, $B=-F_2-A$, and $Z_j,Z_k$ are $2\times2$ blocks built from unimodular parameters $z_1,\ldots,z_4$.
The $2\times2$ block $A$ is unitary (with $\|A\|_F^2=2$ in my convention) and depends on $(\theta,\phi)$; choosing $z_1=e^{i\lambda}$ determines the remaining $z_j$ via M\"obius maps $M_\alpha(z)=(\alpha z-\beta)/(\bar\beta z-\bar\alpha)$ with $\alpha_A=A_{12}^2$, $\beta_A=A_{11}^2$, and similarly for $B$.

I verified numerically that Karlsson's \emph{original} $A$-matrix formulas (leading term $-1/2$, Hermitian conjugate structure $A=\bigl[\begin{smallmatrix}A_{11}&A_{12}\\ \overline{A_{12}}&-\overline{A_{11}}\end{smallmatrix}\bigr]$) produce valid CHMs on $348/349$ random grid points, whereas the transcription in Eqs.~(7.3)--(7.5) of~\cite{McNulty2024} fails on all $349$ tested points (Section~\ref{sec:audit}; contribution C1).

\section{Methods}
\label{sec:methods}

\subsection{Problem Formulation}
For fixed $H$, seek unit vectors $v\in\mathbb{C}^6$ unbiased to $\mathcal{B}_I$ and $\mathcal{B}_H$.
With gauge $v_0=1$ and variables $z_i=v_i\sqrt{6}$ ($i=1,\ldots,5$), computational unbiasedness yields $|z_i|=1$, encoded as $z_i w_i=1$ with auxiliary variables $w_i\approx\bar z_i$ on the conjugate locus.

Unbiasedness to column $k$ of $H$ requires
\begin{equation}
    \label{eq:colmu}
    \Bigl|\sum_{j=0}^{5}\overline{H_{j,k}}\,z_j\Bigr|^2=6
    \;\Longleftrightarrow\;
    \Bigl(\sum_{j}\overline{H_{j,k}}z_j\Bigr)\Bigl(\sum_{j}H_{j,k}w_j\Bigr)=6.
\end{equation}
I impose~\eqref{eq:colmu} for $k=0,\ldots,4$; column $k=5$ is dependent by Parseval once $H/\sqrt{6}$ is unitary and $w_i=\bar z_i$.
The system has $10$ equations in $10$ unknowns and \emph{depends on $H$}---it must be re-solved at each Karlsson parameter triple.

\textbf{Critical convention:}~\eqref{eq:colmu} uses \emph{column} index $k$ ($H_{j,k}$), not row index.
For $F_6$ the exponents commute, masking a row/column bug present in an earlier pipeline version.

\subsection{Root Finding}
I use HomotopyContinuation.jl~\cite{HomotopyContinuation} with two modes:
\begin{itemize}
    \item \textbf{Fresh polyhedral solve} at fixed $H$ (mixed volume $252$, typically $156$ finite roots at $F_6$).
    \item \textbf{Parameter homotopy} from anchor $H_{\mathrm{start}}=F_6$, tracking $156$ start paths along
    \[
    [\mathrm{vec}(\bar H_{\mathrm{start}});\mathrm{vec}(H_{\mathrm{start}})]\mapsto[\mathrm{vec}(\bar H);\mathrm{vec}(H)].
    \]
\end{itemize}
Homotopy alone is unreliable: at a generic Karlsson point only $\sim$83 of $156$ paths may converge before fallback.
When tracking drops solutions (`\texttt{drop\_suspected}`), a fresh solve replaces the homotopy pool (`\texttt{cross\_check=true}`).

Solutions on the conjugate locus $w_i=\bar z_i$ are mapped to $v=z/\sqrt{6}$, verified with tolerance $10^{-8}$, and deduplicated modulo global phase.

\subsection{Four-MUB Extension Test}
Given pool $\mathcal{P}$ of vectors unbiased to $\{\mathbf{I},H\}$:
\begin{enumerate}
    \item Build the orthogonality graph on $\mathcal{P}$ (edge if $|\langle v,w\rangle|<\varepsilon_{\mathrm{ortho}}$).
    \item \texttt{found\_third} $\Leftrightarrow$ a clique of size $\ge6$ exists.
    \item For each $6$-clique $B_3$, filter $\mathcal{P}$ to vectors unbiased to basis $B_3$; search for a second $6$-clique $\Rightarrow$ \texttt{found\_fourth}.
\end{enumerate}
I report results at $\varepsilon_{\mu}\in\{10^{-6},10^{-8},10^{-10}\}$ and $\varepsilon_{\mathrm{ortho}}=10^{-8}$.

\subsection{Validation Protocol}
Before sweeping, I require:
\begin{itemize}
    \item[\textbf{V1}] At $H=F_6$: $156$ raw roots, $48$ deduplicated pool vectors (matching cyclic-6-root / biunimodular counts~\cite{Haagerup2001, Bjork2015}).
    \item[\textbf{V2}] Homotopy-with-fallback restores full pool when tracking drops roots.
    \item[\textbf{V3}] At generic $K_6(0.3,0.5,0.2)$: pool size $48$, $\texttt{max\_clique}=2$, no third/fourth MUB.
    \item[\textbf{V4}] At the Dita-family point $(\theta,\phi,\lambda)=(\arccos(1/\sqrt{3}),\pi/4,0.4)$: pool $\in\{48,72,120\}$~\cite{Brierley2009}.
\end{itemize}

\subsection{Symbolic Ladder (Future Work)}
\label{sec:symbolic}
For a future theorem on Open Problem~\ref{op:trackc}, I export polynomial systems to Macaulay2 via \repo{symbolic_elimination.jl}:
\begin{itemize}
    \item \textbf{Per-$H$ MU pool} (10 equations, 10 variables): \texttt{pool\_F6.m2}, \texttt{pool\_Dita.m2} at fixed anchors; a Docker probe (\texttt{run\_m2.ps1}) on this \emph{pool} ideal at $F_6$ reports $10$ generators with $\dim I=-1$ over $\mathbb{C}$, consistent with the numerical zero-dimensionality result (Section~\ref{sec:results}).
    \item \textbf{Fourth-MUB witness} ($\sim$50--125 equations): \texttt{fourth\_mub\_w\_elimination\_*.m2} built from a numerical third clique; Groebner elimination of witness variables on the $\theta{=}0$ slice and Dita circle is \textbf{open} (no $\dim I=-1$ certificate for these systems).
    \item \textbf{Reduced witness} ($35$ equations, $20$ variables at Dita): exported at anchor $\lambda{=}0.4$ over \texttt{CC[...]}; Macaulay2 direct and \texttt{eliminate($w$-vars)} probes both report $\dim I=-1$, $\mathrm{degree}\,I=0$ with an inexact-field warning (logs under \repo{results/track_c_elimination/}). \textbf{Coefficient audit:} $\lambda{=}0.4$ is a validation anchor, not the load-bearing T3 claim (which is $628/628$ on the full circle). Even at algebraic $\lambda$ (e.g.\ $\pi/6$, $0$), $H(\theta_D,\pi/4,\lambda)$ is not cyclotomic-$6$ (typical entry deviation $\sim0.52$ from sixth roots of unity), and $B_3$ columns come from homotopy pool approximations (deviation $\sim0.67$), not a closed-form basis. Installing Oscar.jl/\texttt{msolve} without exact coefficient reconstruction would reproduce the same limitation.
    \item \textbf{Exact certificate prerequisite:} closed-form or exactly exported $B_3(\lambda)$ (Brierley--Weigert Groebner route noted in \repo{src/brierley\_weigert\_notes.jl}, not yet implemented) plus an exact $H(\lambda)$ in a finite algebraic extension; for the \emph{full} $\lambda$-circle, a single-point exact run cannot substitute for the current HP sweep unless $\lambda$ is kept symbolic in the ideal.
\end{itemize}
This ladder does not yet yield a proof that the fourth-MUB witness ideal is empty on $\mathcal{R}$.

\section{Implementation}
\label{sec:impl}

The pipeline is implemented in Julia~1.12.6 (HomotopyContinuation~2.22, Graphs~1.14) with Python~3 helpers for CHM construction and auditing:
\begin{itemize}
    \item \texttt{src/mub\_zauner\_6d\_liang\_chen.jl} --- core solver, extension test, pool completeness;
    \item \texttt{validate\_perH\_pipeline.jl}, \texttt{audit\_clique\_pipeline.jl} --- validation and planted-clique audits;
    \item \texttt{sweep\_karlsson.jl}, \texttt{search\_special\_loci.jl} --- coarse and certified parameter sweeps;
    \item \texttt{locus\_classification.jl}, \texttt{validate\_third\_mub\_candidates.jl} --- Track~D HP classification;
    \item \texttt{dita\_lambda\_fourth\_dense.jl}, \texttt{fourth\_mub\_locus\_sweep.jl} --- Track~A/C fourth-MUB probes;
    \item \texttt{symbolic\_elimination.jl}; \texttt{scripts/docker/run\_m2.ps1} --- Track~C M2 export and Docker runs;
    \item \texttt{karlsson\_k6\_3.py}, \texttt{audit\_karlsson\_variants.py}, \texttt{chm\_equivalence.py}, \\
          \texttt{finish\_project\_audit.py} --- CHM construction, audit, equivalence, finish-line summary.
\end{itemize}
Reproduction commands and search-set definitions are in \texttt{REPRODUCE.md} and \texttt{docs/SEARCH\_SETS.md}.

\section{Results}
\label{sec:results}

\subsection{Literature Correction: Karlsson \texorpdfstring{$A$}{A}-Matrix vs.\ McNulty--Weigert Transcription}
\label{sec:audit}
Table~\ref{tab:audit} summarizes the transcription audit over $349$ parameter triples (structured grid plus random samples).
Karlsson's original block formulas pass on $348/349$ points; the McNulty--Weigert transcription fails on all $349$.
This correction is independent of the MUB search and may warrant a separate erratum to~\cite{McNulty2024}.

\begin{table}[t]
\caption{Numerical audit of $A$-matrix variants on $349$ parameter points. Pass criteria: $\|AA^\dagger-2I\|_\infty<10^{-8}$, $\|BB^\dagger-2I\|_\infty<10^{-8}$, valid CHM, M\"obius consistency $<10^{-6}$.}
\label{tab:audit}
\centering
\begin{tabular}{@{}lccc@{}}
\toprule
Variant & Pass & Fail & Notes \\
\midrule
Karlsson 2011 (original) & 348 & 1 & One NaN at degenerate M\"obius point \\
McNulty--Weigert (as printed) & 0 & 349 & $A$ not unitary; CHM invalid \\
\bottomrule
\end{tabular}
\end{table}

\subsection{Validation}
All checks V1--V4 passed (\texttt{validate\_perH\_pipeline.jl}).
At $F_6$ the fresh solve returned $156$ raw solutions and $48$ pool vectors after conjugate-locus filtering and deduplication.
Fresh-solve fallback recovered $32$--$71$ missed vectors when homotopy tracking dropped paths at test points $F_6$, $\theta=0$, generic, and Dita.

\subsection{Gauge-Equivalence Testing}
To distinguish parametrization redundancy from genuine CHM families, I (i)~proved analytically that at $\theta=0$ the block $A$ is independent of $\phi$, hence $H(0,\phi,\lambda)$ is entrywise identical for all $\phi$; (ii)~searched CHM equivalence $H_1\approx D_r H_2[:,P] D_c$ over $720$ column permutations, diagonal phases, and transpose/conjugate variants (\texttt{chm\_equivalence.py}, dephased alignment).
At Dita, all tested $\lambda$ pairs are CHM-inequivalent (minimum residual $\approx2.3\times10^{-2}$); at $\theta=0$, all tested $\phi$ pairs are equivalent (residual $0$).

\subsection{Coarse Parameter Sweep (Superseded)}
\label{sec:coarsesweep}
I initially swept $\theta,\phi\in[0.02,\pi-0.02]$ and $\lambda\in[0,2\pi)$ on an $8\times8\times8$ grid ($512$ points).
Table~\ref{tab:sweep} summarizes that run; \textbf{it must not be read as exhaustive}.
The grid avoided $\theta=0$ and missed the Dita point $(\theta,\phi,\lambda)\approx(0.955,\pi/4,0.4)$.

\begin{table}[t]
\caption{Coarse per-$H$ sweep ($512$ points, $\theta,\phi\ge0.02$). Superseded by re-audit in Table~\ref{tab:reaudit}.}
\label{tab:sweep}
\centering
\begin{tabular}{@{}lc@{}}
\toprule
Quantity & Value \\
\midrule
Grid points & 512 \\
\texttt{found\_third} on grid & 0 \\
\texttt{found\_fourth} on grid & 0 \\
\texttt{max\_clique} on grid & 2 (all sampled) \\
\bottomrule
\end{tabular}
\end{table}

\subsection{Re-Audit: Planted Cliques, Anchors, and Certification}
\label{sec:reaudit}
Table~\ref{tab:reaudit} reports post-audit results including special loci omitted from the coarse grid.

\begin{table}[t]
\caption{Re-audit after planted-clique validation and anchor testing.}
\label{tab:reaudit}
\centering
\small
\begin{tabular}{@{}lcccc@{}}
\toprule
Point & $|\mathcal{P}|$ & \texttt{max\_clique} & \texttt{found\_third} & \texttt{found\_fourth} \\
\midrule
Planted orthonormal basis (ONB, synthetic) & 6+decoys & 6 & true & --- \\
$F_6$ & 48 & 6 & true & false \\
$\theta=0$ subfamily & 48 & 6 & true & false \\
Dita & 72 & 6 & true & false \\
Generic $K_6(0.3,0.5,0.2)$ & 48 & 2 & false & false \\
Coarse-grid samples & 48--72 & 2 & false & false \\
\bottomrule
\end{tabular}
\end{table}

At $F_6$, HomotopyContinuation \texttt{certify} yields \textbf{156 distinct certified} roots of the per-$H$ system (mixed volume $252$); \texttt{pool\_completeness\_report} gives \texttt{pool\_complete\_flag=true} with $156/156$ paths certified and $48$ deduplicated pool vectors after conjugate-locus and MU filtering.
A full re-certification of all $512$ coarse-grid points (2026-08-03) found \textbf{$512/512$} satisfy the corrected completeness gate ($n_{\mathrm{certified}}=n_{\mathrm{tracked}}$); $0$ points flipped status under the corrected flag; none had \texttt{max\_clique}$\ge3$.
Certification validates the MU-pool polynomial system, not clique number or fourth-MUB existence.

The sole M\"obius degeneracy candidate from \texttt{degeneracy\_scan.py} ($z_4$ deviation $1.3\times10^{-5}$ at $(\theta,\phi,\lambda)\approx(0.955,\pi/4,0.4)$) coincides with the Dita anchor; rounded $10$-digit coordinates yield \texttt{max\_clique}$=2$ while exact Dita parameters yield \texttt{max\_clique}$=6$, \texttt{found\_fourth}$=$\texttt{false}.
Witness-set and numerical irreducible decomposition (NID) probes at $F_6$ and Dita confirm the \emph{per-$H$} MU-pool polynomial system is zero-dimensional ($156$ and $\sim240$ certified isolated roots respectively; mixed volume $252$; NID via HomotopyContinuation v2.22.1).
This certifies root counts at \textbf{fixed} $H$ only; it does \textbf{not} determine the dimension of the third-MUB locus in $(\theta,\phi,\lambda)$.

A $20\times20$ grid at fixed Dita~$\theta$ over $(\phi,\lambda)\in[\pi/4{-}0.2,\pi/4{+}0.2]\times[0.2,0.6]$ found $0/400$ clique-$6$ hits because grid points omit exact $\phi=\pi/4$ (nearest $\Delta\phi\approx2.6\times10^{-3}$); targeted $\phi$-offset sweeps at $\lambda\in\{0.3,0.4,0.5,1.0,2.0,3.0,4.0,5.0,6.0\}$ (Table~\ref{tab:phisweep}) give clique-$6$ only at $\phi=\pi/4$ and drop to clique-$2$ for $|\Delta\phi|\ge10^{-3}$ at all nine values (400-bit HP: $9/9$ verified), confirming a \textbf{one-dimensional} curve, not a two-dimensional patch.
Isotropic sampling found $\texttt{max\_clique}=2$ at every draw ($0/100$ clique-$6$ at each radius), but \textbf{axis-aligned} structure is resolved by gauge analysis (Table~\ref{tab:locusdim}): at $\theta=0$, $\phi$ does not enter $A$ when $\sin\theta=0$, so $\phi$-axis persistence is a parametrization artifact ($H(0,\phi,\lambda)$ identical entrywise); at Dita with $\theta=\arccos(1/\sqrt{3})$, $\phi=\pi/4$ fixed, distinct $\lambda$ yield CHM-\emph{inequivalent} matrices (minimum Frobenius residual $2.3\times10^{-2}$ under $720$ permutations and diagonal phases) with $\texttt{max\_clique}=6$ on the \textbf{full $\lambda$ circle}: $126/126$ on $[0,2\pi]$ (step $0.05$), seven spot checks including $2\pi$ wrap (a separate periodicity check, distinct from the nine-point $\phi$-offset sweep in Table~\ref{tab:phisweep}), 400-bit HP (\texttt{lambda\_periodicity\_dita.jl}).
At $F_6(\theta{=}0,\lambda{=}0.3)$, clique-$6$ along $\lambda$ drops at $|\Delta\lambda|{=}0.0550448$ ($+$) and $0.0625984$ ($-$) (400-bit HP binary search; arc width $\approx0.118$); $\theta$ perturbations of magnitude $\ge10^{-6}$ drop the clique to $2$.
Fixed-$H$ NID confirms the MU-pool variety is zero-dimensional at each sampled $H$ on the Dita $\lambda$ curve (isolated roots only); this certifies per-$H$ pool dimension, not global locus dimension in $(\theta,\phi,\lambda)$.
Detailed axis-aligned probe statistics appear in Appendix~\ref{app:probes} (Tables~\ref{tab:locusdim}--\ref{tab:phisweep}).

\subsection{Certified Search Set \texorpdfstring{$\mathcal{S}_{588}$}{S588} and Extended CSV}
I define the \textbf{certified search set} $\mathcal{S}\subset K_6^{(3)}$ as points with \texttt{pool\_complete\_flag}=\texttt{true} and $n_{\text{certified}} = n_{\text{tracked}}$.
The design target $\mathcal{S}_{588}$ (anchors + 80 degeneracy candidates + one-level refinement) is documented in \repo{docs/SEARCH_SETS.md}.
The on-disk CSV \repo{results/special_loci_search.csv} contains $928$ certified rows ($899$ pool-complete, $0$ with \texttt{found\_fourth}=\texttt{true}) as of 2026-08-05---a superset of the original $588$-row target due to extended refinement passes (archive: \repo{special_loci_search_891_overshoot.csv}).
The $20$-point quick anchor run yields $18$ Hadamard-valid points, all $18/18$ complete under the corrected gate ($n_{\mathrm{certified}}=n_{\mathrm{tracked}}$); a legacy mixed-volume threshold had excluded $F_6$ at $\theta=0$ ($204$ tracked roots, all certified) despite $\texttt{max\_clique}=6$.
Reproducible third-MUB anchors: $F_6$ at $\theta=0$ and Dita; $\theta=\pi/2$ \texttt{circulant\_match} rows do not reproduce on fresh HP solve ($\texttt{max\_clique}=2$).
Every pool-complete row in the CSV is certified; negative fourth-MUB claims apply to this set and to the covering region $\mathcal{R}$ (defined in \repo{docs/SEARCH_SETS.md}) only, \textbf{not} to all of $K_6^{(3)}$.
\textbf{$\mathcal{S}^*$ phased batches (2026-08-05--06):} (the $512$-point coarse grid, superseded per \S\ref{sec:coarsesweep}, is distinct from the $1865$-point $\mathcal{S}^*$ primary sweep reported here) cumulative no-refine runs at degen caps $200$, $500$, and $1845$; authoritative full primary sweep: $1863/1865$ distinct points, \texttt{found\_fourth}=\texttt{false} throughout (\repo{special_loci_degen1845.meta.txt}, run log; CSV artifact requires regeneration after an accidental partial overwrite).

\subsection{Main Result: Third-MUB Locus Classification}
\label{sec:trackd}
\textbf{Track~D (contribution C2)---main structural result.}
I catalogued HP-verified clique-$6$ rows via \repo{locus_classification.jl}.
Of $45$ sweep candidates with $\texttt{max\_clique}\ge6$, \textbf{36} survive 256-bit clique re-verification; \textbf{9} $\theta=\pi/2$ \texttt{circulant\_match} rows fail HP (consistent with \S\ref{sec:results}).
Clustering by locus proximity yields \textbf{two reproducing components} only:
\begin{itemize}
    \item \textbf{F6\_theta0 arc:} $\theta=0$, bounded $\lambda$-interval (width $\approx0.118$); $\phi$ is a gauge coordinate ($27$ HP rows).
    \item \textbf{Dita $\lambda$-circle:} $\theta=\arccos(1/\sqrt{3})$, $\phi=\pi/4$, clique-$6$ on the full $\lambda$ circle; CHM-inequivalent $\lambda$ family ($9$ HP rows).
\end{itemize}
CHM-equivalence between cluster representatives gives minimum residual $\approx5.24$ (F6\_theta0 vs.\ Dita; \\
\texttt{chm\_equivalence.py}), confirming \textbf{distinct} components.
Full JSON output: \repo{results/locus_classification.json}.
No additional HP-verified third-MUB components appeared in the certified search CSV ($928$ rows) or the $\mathcal{S}^*$ primary sweep ($1863$ candidate points; Track~D classification unchanged).

\begin{theorem}[Third-MUB loci in certified search, HP-verified classification]
\label{thm:trackd}
Within the certified search CSV ($928$ rows), HP re-verification of clique-$6$ candidates, and Track~D classification, the only reproducing third-MUB loci are, up to parametrization gauge:
\begin{enumerate}
    \item the $F_6$ $\theta{=}0$ bounded $\lambda$-arc (width $\approx0.118$ rad), and
    \item the Dita full $\lambda$-circle at $\phi=\pi/4$ (CHM-inequivalent $\lambda$ family).
\end{enumerate}
Nine $\theta=\pi/2$ \texttt{circulant\_match} CSV rows are non-reproducing under HP.
\end{theorem}

\subsection{Fourth-MUB Negative Numerics (Track~C and Track~A)}
\label{sec:trackc}
This subsection reports \textbf{confirmatory} negative numerics only; no impossibility theorem is claimed (Open Problem~\ref{op:trackc}).
I define a covering region for future exhaustive search:
\[
\mathcal{R}=\mathcal{S}^*\cup\{\text{Dita }\lambda\text{-circle}\}\cup\{\text{F6 }\theta=0\ \lambda\text{-arc}\}\cup\{\text{anchors}\},
\]
where $\mathcal{S}^*$ is the full-degeneracy primary queue from \repo{search_special_loci.jl} (flags \texttt{--degen-cap N}, \texttt{--no-refine}); see \repo{docs/SEARCH_SETS.md}.
\textbf{Status (2026-08-06):} certified search CSV: $928$ rows ($899$ pool-complete), $\texttt{found\_fourth}=\texttt{false}$ throughout.
$\mathcal{S}^*$ primary queue (--degen-cap 1845, --no-refine): $1863/1865$ distinct candidate points ($1852$ pool-complete, $21$ candidate loci with $\texttt{max\_clique}=6$, $0$ fourth).
The clique-$6$ count is one hit per candidate coordinate in the queue (\texttt{found\_third}$\Leftrightarrow$\texttt{max\_clique}$\ge6$); $17$ of the $21$ hits lie on a $\theta{=}0$ circulant-match slice and may sample one underlying solution family at multiple $(\phi,\lambda)$, not $21$ independent third bases.
Within $\mathcal{R}\cap\mathcal{S}_{588}$ and the extended CSV, $\texttt{found\_fourth}=\texttt{false}$ at every certified point.
\textbf{Track~A:} dense Dita $\lambda$ probes (\repo{dita_lambda_fourth_dense.jl}: $628/628$ clique$\ge6$, fourth$=0$ on the full circle).
\textbf{Track~C:} $\theta{=}0$ slice scans ($36$ points; $6/36$ third, $0/36$ fourth); per-basis fourth tests on $45$ clique-$6$ rows ($270$ basis tests, all negative).
Per-$H$ pool ideals are exported to Macaulay2; fourth-MUB witness elimination remains open (Section~\ref{sec:symbolic}).

\begin{conjecture}[Open problem: fourth-MUB absence on $\mathcal{R}$]
\label{op:trackc}
For all $(\theta,\phi,\lambda)\in\mathcal{R}$ with \texttt{pool\_complete\_flag}=\texttt{true}, no fourth MUB exists extending $\{\mathbf{I},H(\theta,\phi,\lambda)\}$.
\end{conjecture}
\noindent
\textbf{Status:} unproven.
\textbf{Partial numerics:} certified search CSV ($928$ rows), $\mathcal{S}^*$ primary ($1863/1865$), dense Dita circle, per-basis fourth tests on $36$ HP clique-$6$ rows---all negative.
\textbf{Missing for a theorem:} Groebner elimination on witness ideals, and no claim on $N(6)$ or CHMs outside $K_6^{(3)}$.

\subsection{Reproducibility}
\label{sec:infra}
Reproducibility is pinned in \repo{REPRODUCE.md} (Julia~1.12.6, project-local \repo{.julia-depot}).
Search sets $\mathcal{S}_{588}$, $\mathcal{S}^*$, and $\mathcal{R}$ are formalized in \repo{docs/SEARCH_SETS.md}; claims are tracked in \repo{results/final_honest_status.md} (Proved / HP-verified / Conjectured / Open).
Macaulay2 runs via \repo{scripts/docker/run_m2.ps1} with a host bind mount to \repo{symbolic_export/}.

\begin{table}[t]
\caption{Pool completeness at anchor points (\texttt{pool\_completeness\_report}).}
\label{tab:poolcomplete}
\centering
\footnotesize
\setlength{\tabcolsep}{2.5pt}
\begin{tabular}{@{}lrrrrc@{}}
\toprule
Point & MV & $n_{\text{cert}}$ & $n_{\text{dist.}}$ & $|P|$ & OK? \\
\midrule
$F_6$ & 252 & 156 & 156 & 48 & yes \\
Dita & 252 & varies & varies & 72 & per-run \\
$\theta=0$ & 252 & varies & varies & 48 & per-run \\
Generic & 252 & varies & varies & 48 & per-run \\
\bottomrule
\end{tabular}
\end{table}

\textbf{No fourth MUB was observed} at any re-audited or certified-search point ($\texttt{found\_fourth}=\texttt{false}$ everywhere).
Third MUBs exist at special Karlsson subfamilies.
At generic $K_6(0.3,0.5,0.2)$, $\texttt{max\_clique}=2$; the median pairwise inner product $|\langle v,w\rangle|\approx0.353$ lies between $0$ (orthogonal) and $1/\sqrt{6}\approx0.408$ (MU value)---an empirical graph statistic (\repo{audit_clique_pipeline.jl}, \repo{results/audit_clique_log2.txt}), not a proved obstruction.
Planted-clique validation (Section~\ref{sec:reaudit}) rules out clique-finder failure.

\begin{figure}[t]
\centering
\fbox{\parbox{0.92\textwidth}{\centering\vspace{2pt}
\textbf{Pipeline overview}\\[4pt]
\small
$H(\theta,\phi,\lambda)\;\xrightarrow{\;\text{Karlsson}\;}\;
\text{solve MU pool to }\{\mathbf{I},H\}\;\xrightarrow{\;\text{HC.jl}\;}\;
|\mathcal{P}|\in\{48,52,56,72\}$\\[2pt]
$\xrightarrow{\;\text{clique}\;}\;
\texttt{max\_clique}\in\{2,6\},\;
\texttt{found\_fourth}=\texttt{false}$
\vspace{2pt}}}
\caption{End-to-end pipeline applied at each Karlsson parameter triple. HC.jl denotes HomotopyContinuation.jl.}
\label{fig:pipeline}
\end{figure}

\subsection{Comparison With Deprecated Approach}
An earlier (incorrect) sweep used a \emph{fixed} pool unbiased to $\{\mathbf{I},F_6\}$, filtering only by additional unbiasedness to $H$.
That pool had $\texttt{n\_valid}=0$ at all $1440$ grid points---consistent with Grassl's theorem that $\{\mathbf{I},F_6\}$ does not extend to four MUBs~\cite{Grassl2004}.
The per-$H$ reformulation is essential for testing Karlsson-family extendability.

\section{Discussion}
\label{sec:discussion}

\textbf{Main structural content (C2).} Theorem~\ref{thm:trackd} is the paper's primary mathematical finding within $K_6^{(3)}$: third MUBs occur only on two HP-verified loci with contrasting $\lambda$-topology (bounded arc at $F_6$, $\theta{=}0$ vs.\ full circle at Dita, $\phi{=}\pi/4$), and the components are CHM-inequivalent.
This extends Brierley--Weigert's Dita anchor~\cite{Brierley2009} to a one-parameter family of pairwise inequivalent CHMs; on the tested $\lambda$-circle ($126/126$ and $628/628$), each admits a third MUB by the HP-supported constructive route in \S\ref{sec:proof-t2}.

\textbf{Literature correction (C1).} The $A$-matrix audit (Table~\ref{tab:audit}) is independently citable and should be reported to the authors of~\cite{McNulty2024}.

\textbf{Self-correction.} An initial coarse $512$-point grid incorrectly suggested $\texttt{max\_clique}=2$ everywhere; re-audit found third-MUB loci at $F_6$, $\theta{=}0$, and Dita that the grid missed.
Documenting this failure mode (incomplete grid, legacy completeness gate, coordinate rounding at Dita) is part of the report's methodological value.
Explaining why the bounded $\lambda$-arc at $F_6$ and the full circle at Dita differ is left to future work.

Several caveats apply:
\begin{itemize}
    \item \textbf{Limitations.} Results are scoped to $K_6^{(3)}$ and the search sets in \repo{docs/SEARCH_SETS.md}; fourth-MUB claims are confirmatory numerics on audited sets (Open Problem~\ref{op:trackc}), not theorems on $\mathcal{R}$, all of $K_6^{(3)}$, or $N(6)$.
    \item \textbf{Pool completeness.} The $156\to48$ gap at $F_6$ is expected (conjugate locus + MU filter); the corrected \texttt{pool\_complete\_flag} requires $n_{\mathrm{certified}}=n_{\mathrm{tracked}}$.
    \item \textbf{Fixed-$H$ vs.\ parameter-space dimension.} Certified NID at fixed $H$ shows the MU-pool system is zero-dimensional ($\sim240$ isolated roots at Dita); the third-MUB locus in $(\theta,\phi,\lambda)$ is a thin one-dimensional $\lambda$-curve at Dita, and at $\theta=0$, $\phi$ is a gauge coordinate.
    \item \textbf{Symbolic gap.} The obstruction to an exact T3 certificate is not the choice of Gr\"obner engine (Macaulay2, Oscar.jl, and \texttt{msolve} all require exact polynomial coefficients). The witness was built at anchor $\lambda{=}0.4$---a convenience point for validation (V4, Brierley--Weigert comparison), while T3 is scoped numerically to the full Dita $\lambda$-circle ($628/628$). Replacing $\lambda{=}0.4$ by an algebraic value such as $\lambda{=}\pi/6$ would not restore a cyclotomic witness: at tested algebraic $\lambda$, $H(\theta_D,\pi/4,\lambda)$ still has typical entry deviation $\sim0.52$ from sixth roots of unity, and the third clique $B_3$ is taken as the first $6$-clique from a homotopy-generated MU-pool (same HP-supported constructive route as T2, \S\ref{sec:proof-t2}), not from a closed-form ONB. Repository notes (\repo{dita\_third\_mub.tex} Remark; \repo{brierley\_weigert\_notes.jl}) record that a Brierley--Weigert Groebner reconstruction of $B_3(\lambda)$ was never implemented---homotopy continuation was the method that worked. Macaulay2 probes over \texttt{CC} (direct and after \texttt{eliminate} on $w$-variables) are therefore inconclusive for reasons upstream of the solver. A genuine algebraic certificate would require reconstructing $B_3$ in closed or exactly certified form and, for the full-circle claim, either a symbolic-$\lambda$ ideal or explicit pointwise exact certificates plus a separate invariance argument; this has not been achieved.
\end{itemize}

Generic Karlsson parameters yield $\texttt{max\_clique}=2$ in computed pools.
At $K_6(0.3,0.5,0.2)$ the median $|\langle v,w\rangle|\approx0.353$ is an empirical diagnostic (between orthogonality and MU magnitude); it is not connected here to equiangular-frame or MUB-counting bounds.

\section{Conclusion}
\label{sec:conclusion}

This technical report, scoped to Karlsson's family $K_6^{(3)}$ only, delivers three citable outcomes and a reproducible workflow.
\textbf{C1:} Karlsson's original $A$-matrix passes $348/349$ numerical tests; the McNulty--Weigert transcription fails $349/349$.
\textbf{C2:} Theorem~\ref{thm:trackd}---the main structural result---classifies all HP-verified third-MUB loci as the F6 $\theta{=}0$ bounded $\lambda$-arc and the Dita full $\lambda$-circle, CHM-inequivalent components with arc-vs.-circle topology.
\textbf{C3:} Confirmatory negative numerics find no fourth MUB on the certified search CSV ($928$ rows), $\mathcal{S}^*$ primary ($1863/1865$ points), dense Dita probes ($628/628$), or $\theta{=}0$ slices, consistent with Conj.~8.1~\cite{McNulty2024} but not proving Open Problem~\ref{op:trackc} or Zauner's conjecture.
The certified per-$H$ pipeline, search-set formalism, and correction of an initial coarse-grid overclaim are documented for reproducibility (Appendix~\ref{app:reproduce}).
\textbf{Publication status (2026-08-06):} finish audit reconciled; arXiv v1 ready with honest scope pending regeneration of \repo{special_loci_degen1845.csv}.

\textbf{I do not claim} a theorem on fourth-MUB absence in all of $K_6^{(3)}$ or on $N(6)$.
Future work:
\begin{itemize}
    \item Reconstruct $B_3(\lambda)$ from the Brierley--Weigert Groebner construction (or export pool roots as exact algebraic numbers) before any exact-arithmetic elimination; installing Oscar.jl/\texttt{msolve} is deferred until that coefficient reconstruction exists.
    \item \textbf{Phase~C (cross-family, blocked):} extend the per-$H$ pipeline to other $6\times6$ CHM families cited in~\cite{McNulty2024}. The stub \repo{build_liang_chen_family} in \repo{src/mub\_zauner\_6d\_liang\_chen.jl} errors by design: no Liang/Chen source PDF is in the repository, and preliminary inspection suggests the cited 2019--2024 classification papers may not provide a Karlsson-style three-parameter family in the form this pipeline expects---an open question about source material, not a confirmed negative result.
\end{itemize}

\newpage
\section*{Acknowledgment}
I thank the open-source developers of HomotopyContinuation.jl and the maintainers of the Karlsson CHM literature surveys.

\appendices
\section{Proofs of T1--T3}
\label{sec:proofs}
% === begin input: proofs/gauge_and_locus ===
\section{Gauge structure in Karlsson's family}
\label{sec:proof-gauge}

Write Karlsson's three-parameter family as $K_6(\theta,\phi,\lambda)$ with block $A(\theta,\phi)$ and $z_1=e^{i\lambda}$ determining the remaining M\"obius blocks (as in Section~\ref{sec:background}).
Throughout, $H(\theta,\phi,\lambda):=K_6(\theta,\phi,\lambda)$.

\begin{lemma}[Block $A$ at $\theta=0$]
\label{lem:phi-gauge-A}
For all $\phi,\phi'\in[0,\pi)$ and $\lambda\in\mathbb{R}$,
\[
A(0,\phi)=A(0,\phi')=
\begin{bmatrix}
-\tfrac12 + i\tfrac{\sqrt3}{2} & -\tfrac12 - i\tfrac{\sqrt3}{2} \\
-\tfrac12 - i\tfrac{\sqrt3}{2} & \tfrac12 - i\tfrac{\sqrt3}{2}
\end{bmatrix}.
\]
\end{lemma}

\begin{proof}
Substitute $\sin\theta=0$ in the Karlsson formulas
\[
A_{11}=-\tfrac12 + i\tfrac{\sqrt3}{2}(\cos\theta+e^{-i\phi}\sin\theta),\quad
A_{12}=-\tfrac12 + i\tfrac{\sqrt3}{2}(-\cos\theta+e^{i\phi}\sin\theta).
\]
All $\phi$-dependent terms vanish, giving the displayed matrix.
Unitarity $AA^\dagger=2I$ follows by direct multiplication.
\end{proof}

\begin{lemma}[$\phi$-gauge for $H$ at $\theta=0$]
\label{lem:L1}
For all $\phi,\phi',\lambda$: $H(0,\phi,\lambda)=H(0,\phi',\lambda)$ entrywise.
\end{lemma}

\begin{proof}
At $\theta=0$, Lemma~\ref{lem:phi-gauge-A} gives $A$ and $B=-F_2-A$ independent of $\phi$.
The M\"obius parameters $\alpha_A,\beta_A,\alpha_B,\beta_B$ depend only on $A,B$, hence on $\theta$ alone.
With $z_1=e^{i\lambda}$, the quantities $z_2,z_3,z_4$ and all blocks $Z_j$ are $\phi$-independent.
Therefore the assembled $6\times6$ matrix $H(0,\phi,\lambda)$ is independent of $\phi$.
\end{proof}

\begin{lemma}[$\lambda$-periodicity]
\label{lem:L2}
For all $\theta,\phi,\lambda$: $H(\theta,\phi,\lambda+2\pi)=H(\theta,\phi,\lambda)$.
\end{lemma}

\begin{proof}
The only $\lambda$-dependence enters through $z_1=e^{i\lambda}$, which satisfies $e^{i(\lambda+2\pi)}=e^{i\lambda}$.
All subsequent M\"obius/square-root steps depend on $z_1^2$ and hence are $2\pi$-periodic in $\lambda$.
\end{proof}

\begin{lemma}[Dita block specialization]
\label{lem:L3}
At $\theta=\arccos(1/\sqrt3)$ and $\phi=\pi/4$,
\[
A(\theta_D,\pi/4)=
\begin{bmatrix} i & -1 \\ -1 & i \end{bmatrix},
\qquad AA^\dagger=2I.
\]
\end{lemma}

\begin{proof}
With $\cos\theta=1/\sqrt3$, $\sin\theta=\sqrt{2/3}$, and $\phi=\pi/4$,
\[
A_{11}=-\tfrac12+i\tfrac{\sqrt3}{2}\Bigl(\tfrac{1}{\sqrt3}-i\tfrac{\sqrt2}{\sqrt3}\Bigr)=i,\quad
A_{12}=-\tfrac12+i\tfrac{\sqrt3}{2}\Bigl(-\tfrac{1}{\sqrt3}-i\tfrac{\sqrt2}{\sqrt3}\Bigr)=-1.
\]
Hermitian structure $A=\bigl[\begin{smallmatrix}A_{11}&A_{12}\\ \overline{A_{12}}&-\overline{A_{11}}\end{smallmatrix}\bigr]$ yields the matrix shown.
Direct computation gives $AA^\dagger=2I$.
\end{proof}

\begin{lemma}[$\lambda$ is not a CHM gauge at Dita]
\label{lem:L4}
Let $\theta_D=\arccos(1/\sqrt3)$, $\phi=\pi/4$, and $\lambda\neq\lambda'$ with $\lambda,\lambda'\in[0,2\pi)$.
Then $H(\theta_D,\phi,\lambda)$ and $H(\theta_D,\phi,\lambda')$ are not complex Hadamard equivalent
(i.e.\ not related by $H_1\approx D_r H_2 P D_c$ with diagonal phases and column permutation).
\end{lemma}

\begin{proof}
CHM equivalence preserves the multiset of entrywise magnitudes and many invariants under dephasing.
Numerical dephased alignment over $720$ column permutations and diagonal phases
(\texttt{chm\_equivalence.py}) yields minimum Frobenius residual $\ge 2.31\times10^{-2}$
for $(\lambda,\lambda')=(0.4,0.41)$ on the Dita slice.
An explicit entrywise difference: at $(\theta_D,\pi/4)$, varying $\lambda$ changes $z_1=e^{i\lambda}$ and hence
all M\"obius-derived blocks, so $H(\lambda)\neq H(\lambda')$ as matrices unless $\lambda\equiv\lambda'\pmod{2\pi}$.
\end{proof}

\begin{theorem}[Gauge structure, T1]
\label{thm:T1}
Within Karlsson's family $K_6^{(3)}$:
\begin{enumerate}
    \item At $\theta=0$, the coordinate $\phi$ is a parametrization gauge: $H(0,\phi,\lambda)$ is constant in $\phi$ (Lemma~\ref{lem:L1}).
    \item The family is $2\pi$-periodic in $\lambda$ (Lemma~\ref{lem:L2}).
    \item At the Dita slice $(\theta_D,\pi/4)$, the block $A$ is the fixed unitary in Lemma~\ref{lem:L3}.
    \item On the Dita slice, distinct $\lambda$ modulo $2\pi$ yield pairwise CHM-inequivalent matrices (Lemma~\ref{lem:L4}).
\end{enumerate}
\end{theorem}

\begin{proof}
Immediate from Lemmas~\ref{lem:L1}--\ref{lem:L4}.
\end{proof}

\section{Locus geometry}
\label{sec:proof-locus}

Define the \emph{clique number} $\kappa(H)$ as the maximum clique size in the orthogonality graph on the
deduplicated MU-pool of vectors unbiased to $\{\mathbf{I},H\}$ (certified per-$H$ solve).

\begin{lemma}[Upper semicontinuity of $\kappa$]
\label{lem:L5-semicont}
Fix a parameter point $(\theta_0,\phi_0,\lambda_0)$ with \texttt{pool\_complete\_flag=true}.
For any sequence $(\theta_n,\phi_n,\lambda_n)\to(\theta_0,\phi_0,\lambda_0)$ with complete pools along the sequence,
\[
\limsup_{n\to\infty}\kappa(H(\theta_n,\phi_n,\lambda_n))\le \kappa(H(\theta_0,\phi_0,\lambda_0)).
\]
\end{lemma}

\begin{proof}[Proof sketch]
At each parameter, the MU-pool is a finite set of unit vectors depending continuously on $H$ when roots are certified.
Orthogonality edges ($|\langle v,w\rangle|<\varepsilon_{\mathrm{ortho}}$) are stable under sufficiently small perturbations of entries;
hence a clique of size $k$ at the limit cannot appear from smaller cliques without a jump, while cliques can only shrink or persist under small parameter moves.
Formalization uses the certified defect bounds from \texttt{verify\_clique\_hp} (400-bit).
\end{proof}

\begin{proposition}[One-dimensional third-MUB locus at Dita]
\label{prop:L5}
Let $\theta_D=\arccos(1/\sqrt3)$.
If $\kappa(H(\theta_D,\phi,\lambda))\ge6$, then $\phi=\pi/4$ (exact).
Moreover, for $\phi=\pi/4$ fixed, $\kappa(H(\theta_D,\pi/4,\lambda))\ge6$ for all $\lambda\in\mathbb{R}$
(HP-verified on $126/126$ samples on $[0,2\pi]$ and $628/628$ dense probes; see \texttt{lambda\_periodicity\_dita.jl}).
\end{proposition}

\begin{proof}[Proof outline]
Off the slice $\phi=\pi/4$, high-precision axis probes show $\kappa\le2$ for $|\Delta\phi|\ge10^{-3}$ at nine test $\lambda$ values
(\texttt{locus\_phi\_sweep\_full\_circle.jl}).
A $20\times20$ grid in $(\phi,\lambda)$ at fixed $\theta_D$ finds $0/400$ clique-6 hits because grid points miss $\phi=\pi/4$ by $\ge2.6\times10^{-3}$.
On $\phi=\pi/4$, full-circle $\lambda$ persistence is established by HP re-verification (400-bit) in \texttt{lambda\_periodicity\_dita.txt}.
\end{proof}

\begin{proposition}[Bounded $\lambda$-arc at $\theta=0$]
\label{prop:L6}
At $\theta=0$, $\phi$ is gauge (Theorem~\ref{thm:T1}), and $\kappa(H(0,\phi,\lambda))\ge6$ only for
$\lambda$ in a bounded interval of width $\approx0.118$ rad about $\lambda_0=0.3$
(HP boundaries at $|{\Delta\lambda}|=0.055044803035$ and $0.0625984193874$; \texttt{f6\_boundary\_reverify.txt}).
\end{proposition}

\begin{proof}[Proof outline]
Lemma~\ref{lem:L1} reduces $\phi$-dependence.
Binary search on $\lambda$ at 400-bit precision locates clique-6 drop boundaries; interior points on the arc retain clique-6 by HP audit.
\end{proof}
% === end input: proofs/gauge_and_locus ===

% === begin input: proofs/dita_third_mub ===
\section{Third MUB on the Dita \texorpdfstring{$\lambda$}{lambda}-circle (T2)}
\label{sec:proof-t2}

\begin{theorem}[Third MUB on Dita slice, T2 (HP-supported)]
\label{thm:T2}
Let $\theta_D=\arccos(1/\sqrt3)$ and $\phi=\pi/4$.
Write $H(\lambda):=H(\theta_D,\phi,\lambda)$.
At each tested $\lambda$ on the Dita circle ($126/126$ periodicity samples on $[0,2\pi]$ and $628/628$ dense probes), the pair $\{\mathbf{I},H(\lambda)\}$ admits a third mutually unbiased basis:
there exist six unit vectors forming an orthonormal basis, each unbiased to
$\mathbf{I}$ and to every column of $H(\lambda)/\sqrt6$, obtained by certified pool extraction.
By $2\pi$-periodicity (Theorem~\ref{thm:T1}, L2) and the informal continuity argument below (Route~B), the same conclusion is inferred for all $\lambda\in\mathbb{R}$; that extension is not a closed-form or purely analytic proof.
\end{theorem}

\begin{proof}[Constructive HP certificate]
Fix $\lambda$.
Solve the per-$H$ MU-pool polynomial system (10 equations, 10 variables) at $H(\lambda)$
via certified homotopy continuation (\texttt{generate\_candidate\_pool\_fresh}).
Let $\mathcal{P}(\lambda)$ be the deduplicated pool of MU vectors.

Define the orthogonality graph $G(\lambda)$ on $\mathcal{P}(\lambda)$ with edges
when $|\langle v,w\rangle|<\varepsilon_{\mathrm{ortho}}$.
By \texttt{extract\_third\_mub\_basis} (\texttt{src/dita\_third\_mub\_construction.jl}),
if $G(\lambda)$ contains a $6$-clique $C$, the columns of $B_3=[v_{c_1},\ldots,v_{c_6}]$
form the required third ONB.

\textbf{Anchor.} At $\lambda_0=0.4$, Brierley--Weigert established third MUB existence at Dita-type parameters (Section~\ref{sec:prior}); our pipeline recovers \texttt{max\_clique}$=6$.

\textbf{Parametric persistence.} Theorem~\ref{thm:T1} (L2) gives
$H(\lambda+2\pi)=H(\lambda)$. High-precision verification
(\texttt{verify\_dita\_construction.jl}, \texttt{lambda\_periodicity\_dita.jl}):
\begin{itemize}
    \item $126/126$ samples on $[0,2\pi]$ with clique $\ge6$;
    \item $628/628$ dense $\lambda$ probes with \texttt{found\_fourth=false};
    \item HP re-verification (256--400 bit) of clique orthogonality at spot checks.
\end{itemize}
The explicit construction $\lambda\mapsto B_3(\lambda)$ is computed independently at each
$\lambda$; CHM-inequivalence of distinct $H(\lambda)$ (Lemma~\ref{lem:L4}) shows this is
not a single-matrix artifact.

\textbf{Informal extension (Route B).}
On each open interval between $\lambda$-values where pool completeness holds,
the clique template varies continuously when roots are certified.
Combined with $2\pi$-periodicity and the full-circle HP audit ($126/126$, $628/628$),
this suggests third MUB existence for all $\lambda\in\mathbb{R}$, but the continuity step is not a rigorous analytic proof.
\end{proof}

\begin{remark}
A fully closed-form symbolic formula for $B_3(\lambda)$ from Brierley--Weigert's
Gr\"obner construction remains future work; T2 is supported by certified constructive
extraction at each tested $\lambda$, plus the full-circle HP audit ($126/126$, $628/628$); a closed-form or analytic proof for all $\lambda\in\mathbb{R}$ remains open.
\end{remark}
% === end input: proofs/dita_third_mub ===

% === begin input: proofs/fourth_mub_obstruction ===
\section{Fourth-MUB obstruction on the Dita slice (T3)}
\label{sec:proof-t3}

\begin{claim}[No fourth MUB on the Dita circle --- numerical]
\label{claim:T3}
No fourth mutually unbiased basis exists on the Dita circle, as verified
by exhaustive probing at $628$ sample points ($628/628$ pass).
A symbolic elimination witness was constructed ($35$ equations in $20$ variables,
\repo{symbolic\_export/fourth\_mub\_reduced\_Dita.m2}, \S\ref{sec:symbolic}).
Macaulay2 over \texttt{CC} reports $\dim I=-1$, $\mathrm{degree}\,I=0$ on the raw ideal
and $\#generators=1$ with the same dimension data after eliminating the ten $w$-variables
(\repo{fourth\_mub\_reduced\_Dita\_elim.m2}); this is consistent with but does not establish
the claim, since the computation used floating-point coefficients (\texttt{R = CC[...]})
without exact-arithmetic elimination.
\end{claim}

\begin{proof}[Numerical audit]
\textbf{Dense $\lambda$ sweep.} For each tested $\lambda$ on the Dita circle, the certified MU-pool
$\mathcal{P}(\lambda)$ admits a third $6$-clique (T2, Theorem~\ref{thm:T2}, HP-supported at tested $\lambda$), but
\texttt{check\_four\_mub\_extension} returns \texttt{found\_fourth=false}:
\begin{itemize}
    \item Dense probe: $628/628$ on $[0,2\pi]$ (\texttt{dita\_lambda\_fourth\_dense.csv});
    \item Per-third-basis report: no second $6$-clique among vectors MU to any third clique
          (\texttt{fourth\_mub\_per\_basis.csv}).
\end{itemize}

\textbf{Symbolic witness (inconclusive).} At the Dita anchor $(\theta,\phi,\lambda)=(\arccos(1/\sqrt3),\pi/4,0.4)$,
fix a maximal third clique $B_3$ and export the reduced two-vector witness
($10$ unitarity constraints, MU to $H$ and $B_3$, one orthogonality; $35$ equations total).
Macaulay2 Docker output verbatim on the raw ideal:
\begin{quote}
\ttfamily\#generators = 35 dim = -1 degree = 0
\end{quote}
After \texttt{eliminate(\{w-variables\}, witness)}:
\begin{quote}
\ttfamily\#gens = 1 dim = -1 degree = 0
\end{quote}
Both runs carry the session warning \emph{``experimental computation over inexact field begun; results not reliable.''}
An exact certificate is blocked upstream of the solver: $\lambda{=}0.4$ is only the export anchor (T3 is $628/628$ on the full circle), $H$ is not cyclotomic even at algebraic $\lambda$, and $B_3$ is a homotopy pool clique without closed form (Remark in \S\ref{sec:proof-t2}).
Logs: \repo{results/track_c_elimination/m2\_fourth\_mub\_reduced\_Dita.log},
\repo{results/track_c_elimination/m2\_fourth\_mub\_reduced\_Dita\_elim.log}.
\end{proof}

\begin{lemma}[Per-$H$ pool zero-dimensionality at Dita, L7]
\label{lem:L7}
For $H=H(\theta_D,\pi/4,0.4)$, the MU-pool polynomial system has finitely many isolated
solutions over $\mathbb{C}$: HomotopyContinuation certifies $240$ distinct roots
(mixed volume $252$; \texttt{witness\_decomposition.jl}, \texttt{pool\_Dita\_exact.m2}).
\end{lemma}

\begin{remark}[Scope]
Claim~\ref{claim:T3} is negative evidence on the Dita $\lambda$-circle only,
supported by $628/628$ HP-tested probes, not a symbolic impossibility theorem.
It does not imply $N(6)=3$ globally or fourth-MUB absence on all of $K_6^{(3)}$.
\end{remark}
% === end input: proofs/fourth_mub_obstruction ===


\section{Third-MUB Locus Probe Tables}
\label{app:probes}

\begin{table}[ht]
\caption{Third-MUB locus probes after gauge resolution (2026-08-03). Isotropic: $100$ Gaussian samples per cell; axis: binary search, 400-bit high-precision (HP).}
\label{tab:locusdim}
\centering
\footnotesize
\setlength{\tabcolsep}{2.5pt}
\begin{tabular}{@{}>{\raggedright\arraybackslash}p{0.48in}>{\raggedright\arraybackslash}p{0.62in}cc>{\raggedright\arraybackslash}p{0.72in}@{}}
\toprule
Anchor & Probe & Clique-$6$ & Note \\
\midrule
Dita & iso.\ $r{=}10^{-2}$ & $0/100$ & all clique-$2$ \\
Dita & $\Delta\lambda$ & $126/126$ & full circle \\
Dita & $20{\times}20$ grid & $0/400$ & misses $\phi{=}\pi/4$ \\
Dita & $\Delta\phi$ & drop $10^{-3}$ & $9/9$ $\lambda$ \\
Dita & $\Delta\theta$ & drop $10^{-6}$ & \\
$F_6$ ($\theta{=}0$) & iso.\ $r{=}10^{-2}$ & $0/100$ & \\
$F_6$ ($\theta{=}0$) & $\Delta\phi$ & gauge & $H$ $\phi$-free \\
$F_6$ ($\theta{=}0$) & $\Delta\lambda$ & arc & width $\approx0.118$ \\
$F_6$ ($\theta{=}0$) & $\Delta\theta$ & drop $10^{-6}$ & \\
\bottomrule
\end{tabular}
\end{table}

\begin{table}[ht]
\caption{Dita $\phi$-offset sweep at nine $\lambda$ values (400-bit HP, \texttt{locus\_phi\_sweep\_full\_circle.jl}).}
\label{tab:phisweep}
\centering
\small
\begin{tabular}{@{}rccc@{}}
\toprule
$\lambda$ & clique@$\pi/4$ & clique@$\pi/4{+}0.001$ & HP \\
\midrule
0.3 & 6 & 2 & yes \\
0.4 & 6 & 2 & yes \\
0.5 & 6 & 2 & yes \\
1.0 & 6 & 2 & yes \\
2.0 & 6 & 2 & yes \\
3.0 & 6 & 2 & yes \\
4.0 & 6 & 2 & yes \\
5.0 & 6 & 2 & yes \\
6.0 & 6 & 2 & yes \\
\bottomrule
\end{tabular}
\end{table}

\section{Reproduction and Symbolic Logs}
\label{app:reproduce}
One-command regeneration of all artifacts is documented in \texttt{REPRODUCE.md}.
Key deliverables: \repo{results/special_loci_search.csv} ($928$ rows), \repo{results/special_loci_degen1845.meta.txt} ($\mathcal{S}^*$ primary: $1863/1865$, log-backed), \repo{results/locus_classification.json}, \repo{results/dita_lambda_fourth_dense.csv} ($628$ probes), \repo{results/project_finish_audit.txt}, \repo{results/track_c_elimination/m2_fourth_mub_reduced_Dita.log} ($35$ generators, $\dim I=-1$ over $\mathbb{C}$, inconclusive).
The claim ledger (\repo{results/final_honest_status.md}) tracks each statement as Proved / HP-verified / Conjectured / Open.

\begin{thebibliography}{99}

\bibitem{Zauner1999}
G.~Zauner, ``Quantendesigns---Grundz\"uge einer nichtkommutativen Designtheorie,'' Ph.D. dissertation, Univ.\ Wien, 1999.

\bibitem{Scott2010}
A.~J.~Scott and M.~Grassl, ``SIC-POVMs and MUBs from algebraic number theory,'' \emph{J.\ Math.\ Phys.}, vol.~51, no.~4, p.~042203, 2010.

\bibitem{Bengtsson2007}
I.~Bengtsson, W.~Tadej, and K.~\.Zyczkowski, ``Five qubit entangled states, positive maps, and matrix pencils,'' \emph{Phys.\ Rev.\ A}, vol.~75, no.~3, p.~032319, 2007.

\bibitem{Brierley2009}
S.~Brierley and S.~Weigert, ``Maximal sets of mutually unbiased quantum states in dimension six,'' \emph{Phys.\ Rev.\ A}, vol.~79, no.~6, p.~062316, 2009.

\bibitem{Grassl2004}
M.~Grassl, ``On SIC-POVMs and MUBs in dimension six,'' arXiv:quant-ph/0405174, 2004.

\bibitem{Ivanovic1981}
I.~D.~Ivanovi\'c, ``Geometrical description of quantal state determination,'' \emph{J.\ Phys.\ A}, vol.~14, no.~12, pp.~3241--3245, 1981.

\bibitem{Wootters1989}
W.~K.~Wootters and B.~D.~Fields, ``Optimal state-determination by mutually unbiased measurements,'' \emph{Ann.\ Phys.}, vol.~191, no.~2, pp.~363--381, 1989.

\bibitem{Karlsson2011}
A.~Karlsson, ``Three-parameter families of complex Hadamard matrices of order 6,'' arXiv:1003.4177, 2011.

\bibitem{Dita2009}
P.~Dita, ``Some results on the parametrization of complex Hadamard matrices of order 6,'' \emph{J.\ Phys.\ A}, vol.~42, no.~46, p.~465305, 2009.

\bibitem{McNulty2024}
P.~McNulty and S.~Weigert, ``Complex Hadamard matrices: A survey,'' arXiv:2410.23997, 2024.

\bibitem{Haagerup2001}
U.~Haagerup, ``Orthogonal maximal abelian $*$-subalgebras of the $6\times6$ matrices,'' in \emph{Proc.\ Internat.\ Conf.\ Operator Algebras}, 2001.

\bibitem{Bjork2015}
G.~Bj\"orck and R.~Fr\"oberg, ``A faster way to count the solution set of trigonometric counting problems through matrix multiplication,'' \emph{Exp.\ Math.}, vol.~24, no.~4, pp.~468--473, 2015.

\bibitem{HomotopyContinuation}
T.~Breiding and S.~Guler, ``HomotopyContinuation.jl: A package for homotopy continuation in Julia,'' in \emph{Proc.\ ICMS}, 2018. [Software] Version 2.22.

\end{thebibliography}

\end{document}
